% !TEX program = xelatex
\documentclass[12pt, a4paper]{article}
\usepackage[utf8]{inputenc}
\usepackage[T1]{fontenc}
\usepackage[spanish]{babel}
\usepackage{xeCJK}
\xeCJKsetup{AutoFakeBold=true}
\setCJKmainfont{Noto Serif SC}
\setCJKsansfont{Noto Serif SC}
\setCJKmonofont{Noto Serif SC}
\usepackage{geometry}
\geometry{margin=2.5cm}
\usepackage{titlesec}
\usepackage{enumitem}
\usepackage{booktabs}
\usepackage{array}
\usepackage{xcolor}
\usepackage{mdframed}
\usepackage{parskip}
\usepackage{fancyhdr}
\usepackage{multicol}

\definecolor{azul}{RGB}{30, 90, 160}
\definecolor{grisclaro}{RGB}{245, 245, 245}
\definecolor{azulclaro}{RGB}{220, 235, 255}
\definecolor{verdesuave}{RGB}{210, 240, 215}
\definecolor{verdeborde}{RGB}{50, 140, 60}
\definecolor{rojosuave}{RGB}{200, 50, 50}

\titleformat{\section}{\large\bfseries\color{azul}}{}{0em}{}[\titlerule]
\titleformat{\subsection}{\normalsize\bfseries\color{azul!80}}{}{0em}{}

\setlength{\headheight}{14.5pt}
\addtolength{\topmargin}{-2.5pt}
\pagestyle{fancy}
\fancyhf{}
\rhead{\textcolor{azul}{\textbf{Chino Mandarín — Unidad 3}}}
\lhead{\textcolor{gray}{Guía 2}}
\rfoot{\thepage}
\lfoot{\textcolor{gray}{\small Repaso integrado}}

\newmdenv[backgroundcolor=azulclaro, linecolor=azul, linewidth=1pt, innerleftmargin=10pt, innerrightmargin=10pt, innertopmargin=8pt, innerbottommargin=8pt]{cuadroazul}
\newmdenv[backgroundcolor=grisclaro, linecolor=gray!40, linewidth=0.5pt, innerleftmargin=10pt, innerrightmargin=10pt, innertopmargin=8pt, innerbottommargin=8pt]{cuadrogris}
\newmdenv[backgroundcolor=verdesuave, linecolor=verdeborde, linewidth=1pt, innerleftmargin=10pt, innerrightmargin=10pt, innertopmargin=8pt, innerbottommargin=8pt]{cuadroverde}

\begin{document}

\begin{center}
{\LARGE \textbf{\textcolor{azul}{Guía de Estudio 2}}}\\[4pt]
{\large Repaso integrado}\\[2pt]
{\normalsize Tiempo estimado: 3--4 horas}
\end{center}

\vspace{0.5em}
\hrule height 2pt
\vspace{1em}

\begin{cuadroverde}
\textbf{Objetivo de hoy:} No aprender contenido nuevo, sino consolidar lo de ayer, automatizar la producción de oraciones y preparar las preguntas que podría hacerte el profesor mañana. Al final de esta guía debes poder responder cualquier pregunta sobre fechas, deseos y compras sin pausar a pensar la gramática.
\end{cuadroverde}

%---------------------------------------------------------
\section{Repaso rápido — Bloque 0 (15 min, sin cuaderno)}
%---------------------------------------------------------

Antes de escribir nada, haz este repaso mental. Para cada punto, recita en voz alta sin mirar nada:

\begin{enumerate}[leftmargin=1.5cm]
\item Los 7 días de la semana en chino, en orden.
\item Los 12 meses del año (solo di el número + 月).
\item Las 3 partes del día que aprendiste (mañana, tarde, noche).
\item Las palabras hoy, ayer y mañana.
\item La estructura para preguntar fecha: ¿cómo preguntas qué día es? ¿cómo preguntas qué día de la semana?
\item La estructura con 想: ¿cómo dices que quieres hacer algo?
\item ¿Cómo preguntas el precio de algo?
\end{enumerate}

Si en algún punto dudas más de 5 segundos, vuelve a la guía del día anterior y repasa esa sección antes de continuar.

%---------------------------------------------------------
\section{Bloque 1 — Escritura sin referencia (aprox. 45 min)}
%---------------------------------------------------------

\subsection{Actividad 1 — Dictado de caracteres: todo el vocabulario nuevo}

Escribe en tu cuaderno los caracteres que corresponden a cada palabra o frase. Sin mirar nada. Al terminar, verifica con el material del día anterior y anota tus errores. Los caracteres con error los escribes 5 veces más.

\begin{enumerate}[leftmargin=1.5cm, itemsep=2pt]
\item hoy
\item ayer
\item mañana
\item escuela / colegio
\item leer libros
\item mañana (parte del día)
\item tarde (parte del día)
\item noche
\item querer / desear (verbo)
\item beber
\item comer
\item té
\item arroz cocido
\item comida china
\item tienda
\item comprar
\item vaso / taza
\item edredón / manta
\item ¿cuánto cuesta?
\item yuan (forma coloquial)
\item lunes
\item miércoles
\item viernes
\item domingo
\item martes
\item jueves
\item sábado
\end{enumerate}

\subsection{Actividad 2 — Escritura de los diálogos completos}

Escribe de memoria, sin referencia, los seis diálogos de las lecciones 7 y 8. No línea por línea: escribe el diálogo completo, luego verifica. Cada línea con error: escríbela 5 veces. Objetivo: cero errores en los caracteres que ya practicaste ayer.

\newpage

%---------------------------------------------------------
\section{Bloque 2 — Gramática productiva (aprox. 60 min)}
%---------------------------------------------------------

\subsection{Actividad 3 — Paradigma completo}

Para cada verbo de la lista, construye en tu cuaderno las 5 oraciones del paradigma. Usa sujeto real (yo, tú, él, mi mamá, mi amigo, etc.) y complementos reales.

\begin{cuadrogris}
\textbf{Las 5 formas del paradigma para cada verbo:}\\
1. Oración afirmativa \quad 2. Oración negativa (con 不)\\
3. Pregunta con 吗 \quad 4. Pregunta con pronombre interrogativo (什么 / 哪 / 几)\\
5. Respuesta a esa pregunta
\end{cuadrogris}

\noindent Verbos a trabajar:

\begin{enumerate}[leftmargin=1.5cm]
\item 喝 (beber)
\item 吃 (comer)
\item 去 + lugar (ir a algún lugar)
\item 买 (comprar)
\item 看书 (leer libros)
\end{enumerate}

\textit{Nota:} Con 想, puedes (y debes) combinar: 想喝, 想吃, 想去, 想买. Úsalo en al menos 3 de los 5 verbos.

\subsection{Actividad 4 — Combinaciones con tiempo}

Escribe en tu cuaderno 10 oraciones originales. Cada oración debe incluir obligatoriamente: una expresión de tiempo (día + parte del día cuando sea posible) + sujeto + verbo + objeto. Ninguna puede repetir la misma combinación de tiempo + verbo.

\begin{cuadrogris}
\textbf{Ejemplo del tipo de oración esperada:}\\
星期四下午我想去朋友家吃中国菜。\\
\textit{El jueves por la tarde quiero ir a casa de mi amigo a comer comida china.}
\end{cuadrogris}

\subsection{Actividad 5 — Traducción activa al chino}

Traduce las siguientes oraciones al chino en tu cuaderno. Sin mirar referencia. Son oraciones que combinan todo el contenido visto hasta ahora.

\begin{enumerate}[leftmargin=1.5cm]
\item Hoy es lunes 23 de junio.
\item Ayer fue domingo y mañana es martes.
\item Esta tarde quiero ir a la escuela a leer libros.
\item ¿Qué quieres comer esta noche?
\item Mi compañero de clase quiere comprar un vaso.
\item ¿Cuánto cuesta ese edredón?
\item El sábado por la mañana no quiero ir a la escuela.
\item ¿El profesor es chino o estadounidense?
\item Mi mamá no habla chino, pero quiere comer comida china.
\item Esta tarde quiero ir a la tienda, ¿quieres venir?\quad \textit{Nota: "¿quieres venir?"}\textit{ = }{你想去吗？}
\item Quiero comprar un vaso. ¿Cuánto cuesta este?
\item ¿Qué día de la semana es tu clase de chino?
\end{enumerate}

%---------------------------------------------------------
\section{Bloque 3 — Práctica de preguntas del profesor (aprox. 45 min)}
%---------------------------------------------------------

\subsection{Actividad 6 — Simulación de preguntas en clase}

El profesor te puede hacer cualquiera de estas preguntas. Escribe la respuesta en tu cuaderno \textbf{primero}, luego di la respuesta en voz alta con fluidez. El objetivo es que la respuesta oral salga sin leer lo que escribiste.

\begin{enumerate}[leftmargin=1.5cm]
\item 今天几号？
\item 今天星期几？
\item 昨天是几月几号？
\item 明天是几月几号，星期几？
\item 你的汉语课是星期几？
\item 明天你去学校吗？
\item 你去学校做什么？
\item 今天下午你想做什么？
\item 今天晚上你想吃什么？
\item 星期六你想做什么？
\item 星期六上午你想去哪里？\quad (哪里 = dónde, nǎlǐ — vocabulario de apoyo)
\item 你想买什么？
\item 你想喝什么？
\item 你想吃中国菜吗？
\item 你会做中国菜吗？\quad (repaso examen 2)
\item 你家有几口人？\quad (repaso examen 2)
\item 你今年多大了？\quad (repaso examen 2)
\item 谁是你的汉语老师？\quad (repaso examen 2)
\end{enumerate}

\subsection{Actividad 7 — Preguntas cruzadas con precios}

Escribe un diálogo propio en tu cuaderno donde preguntas por el precio de al menos 3 objetos distintos. Puedes inventar los precios. El diálogo debe tener mínimo 8 líneas e incluir 这个 y 那个.

%---------------------------------------------------------
\section{Bloque 4 — Integración total (aprox. 45 min)}
%---------------------------------------------------------

\subsection{Actividad 8 — Carta breve en chino}

Escribe en tu cuaderno un párrafo de 8 a 10 oraciones en chino hablando sobre tu semana. Debe incluir:

\begin{itemize}[leftmargin=1.5cm]
\item Qué día es hoy y qué hiciste ayer.
\item Qué quieres hacer esta tarde y esta noche.
\item Qué harás mañana y el martes (tu clase de chino).
\item Qué quieres comer o beber.
\item Algo sobre tu familia (vocabulario examen 2).
\end{itemize}

No busques vocabulario nuevo. Si no sabes cómo decir algo, usa lo que sí sabes para rodear la idea.

\subsection{Actividad 9 — Repaso final de caracteres críticos del examen}

Escribe una vez cada uno de estos caracteres sin referencia. Son los del examen 1 y 2 que debes tener ya automatizados:

\begin{cuadrogris}
你好 · 谢谢 · 对不起 · 不客气 · 再见 · 我们 · 中国 · 名字 · 老师 · 学生\\
朋友 · 同学 · 妈妈 · 爸爸 · 哥哥 · 女儿 · 家 · 说汉语 · 中国菜 · 会
\end{cuadrogris}

Si tienes error en alguno, ese carácter va a tu lista de pendientes para repasar antes de cada clase.

\subsection{Actividad 10 — Preguntas de reordenamiento final}

Reordena y escribe la traducción al español:

\begin{enumerate}[leftmargin=1.5cm]
\item 多少钱 / 这个 / 被子
\item 晚上 / 吃 / 想 / 今天 / 我 / 中国菜
\item 去 / 下午 / 商店 / 你 / 吗 / 今天
\item 星期六 / 做 / 你 / 上午 / 什么 / 想
\item 买 / 杯子 / 想 / 我的 / 朋友 / 一个
\item 学校 / 我 / 看书 / 去 / 明天 / 想
\item 几号 / 今天 / 是 / 几月
\item 她 / 会 / 不 / 汉语 / 说 / 妈妈 / 的 / 我
\end{enumerate}

\vspace{1em}
\hrule
\vspace{0.5em}

\begin{cuadroverde}
\textbf{Antes de dormir (10 min):} Escribe de memoria los 7 días de la semana, los 3 momentos del día y las palabras hoy/ayer/mañana. Si salen todos sin error, estás listo para la clase del martes. Si hay algún error, ese punto lo repasas 5 minutos mañana antes de salir.
\end{cuadroverde}

\newpage
%---------------------------------------------------------
\section*{\textcolor{rojosuave}{Respuestas}}
%---------------------------------------------------------

\textbf{Actividad 1 — Dictado de caracteres:}\\
\begin{tabular}{lll}
1. 今天 & 2. 昨天 & 3. 明天 \\
4. 学校 & 5. 看书 & 6. 上午 \\
7. 下午 & 8. 晚上 & 9. 想 \\
10. 喝 & 11. 吃 & 12. 茶 \\
13. 米饭 & 14. 中国菜 & 15. 商店 \\
16. 买 & 17. 杯子 & 18. 被子 \\
19. 多少钱 & 20. 块 & 21. 星期一 \\
22. 星期三 & 23. 星期五 & 24. 星期日 \\
25. 星期二 & 26. 星期四 & 27. 星期六 \\
\end{tabular}

\vspace{1em}
\textbf{Actividad 5 — Traducción al chino:}
\begin{enumerate}[leftmargin=1.5cm]
\item 今天是六月二十三号，星期一。
\item 昨天是星期日，明天是星期二。
\item 今天下午我想去学校看书。
\item 今天晚上你想吃什么？
\item 我的同学想买一个杯子。
\item 那个被子多少钱？
\item 星期六上午我不想去学校。
\item 老师是中国人还是美国人？\quad (还是 = o [en pregunta], háishi — vocabulario de apoyo)
\item 我妈妈不会说汉语，但是她想吃中国菜。
\item 今天下午我想去商店，你想去吗？
\item 我想买一个杯子，这个多少钱？
\item 你的汉语课是星期几？
\end{enumerate}

\vspace{0.5em}
\textbf{Actividad 6 — Preguntas de clase:}
\begin{enumerate}[leftmargin=1.5cm]
\item 今天六月二十三号。
\item 今天星期一。
\item 昨天是六月二十二号。
\item 明天是六月二十四号，星期二。
\item 我的汉语课是星期二和星期四。
\item 明天我去学校。
\item 我去学校学习汉语。
\item 今天下午我想学习汉语。 (o según tu caso)
\item 今天晚上我想吃米饭。 (o según tu caso)
\item 星期六我想去朋友家。 (o según tu caso)
\item 星期六上午我想去商店。 (o según tu caso)
\item 我想买一个杯子。 (o según tu caso)
\item 我想喝茶。 (o según tu caso)
\item 我想吃中国菜。 / 我不想吃中国菜。
\item 我不会做中国菜。 / 我会做中国菜。
\item 我家有 \_\_ 口人。 (número real)
\item 我今年 \_\_ 岁了。 (tu edad)
\item 我的汉语老师是李月老师。
\end{enumerate}

\vspace{0.5em}
\textbf{Actividad 10 — Reordenamiento:}
\begin{enumerate}[leftmargin=1.5cm]
\item 这个被子多少钱？\quad \textit{¿Cuánto cuesta este edredón?}
\item 今天晚上我想吃中国菜。\quad \textit{Esta noche quiero comer comida china.}
\item 今天下午你去商店吗？\quad \textit{¿Esta tarde vas a la tienda?}
\item 星期六上午你想做什么？\quad \textit{¿Qué quieres hacer el sábado por la mañana?}
\item 我的朋友想买一个杯子。\quad \textit{Mi amigo quiere comprar un vaso.}
\item 明天我想去学校看书。\quad \textit{Mañana quiero ir a la escuela a leer libros.}
\item 今天是几月几号？\quad \textit{¿Qué fecha es hoy?}
\item 我妈妈不会说汉语。\quad \textit{Mi mamá no sabe hablar chino.}
\end{enumerate}

\end{document}